% Exported from Math++. Compile with pdflatex.
\documentclass[11pt,a4paper]{article}
\usepackage[T1]{fontenc}
\usepackage[utf8]{inputenc}
\usepackage{lmodern}
\usepackage{amsmath,amssymb,amsthm,mathtools}
\usepackage{graphicx,booktabs,xcolor}
\usepackage[hidelinks]{hyperref}
\usepackage[margin=25mm]{geometry}
\theoremstyle{plain}
\newtheorem{theorem}{Theorem}
\newtheorem{lemma}[theorem]{Lemma}
\newtheorem{corollary}[theorem]{Corollary}
\newtheorem{proposition}[theorem]{Proposition}
\theoremstyle{definition}
\newtheorem{definition}[theorem]{Definition}
\newtheorem{example}[theorem]{Example}
\theoremstyle{remark}
\newtheorem*{remark}{Remark}
\setlength{\parindent}{0pt}\setlength{\parskip}{0.6em}
\begin{document}
\section*{Derivatives}

The derivative of a function, written as you would write it.

\begin{equation*}
  \frac{d}{dx} x^3 \sin x = x^{3} \cos\left(x\right) + 3 x^{2} \sin\left(x\right)
\end{equation*}
\begin{equation*}
  \frac{d}{dx} e^{x^2} = 2 e^{x^{2}} x
\end{equation*}
\begin{equation*}
  \frac{d}{dx} \ln(\cos x) = -\frac{\sin\left(x\right)}{\cos\left(x\right)}
\end{equation*}
Define a function, then differentiate it:

\begin{equation*}
  f(x) \coloneqq  x^2 \ln x
\end{equation*}
\begin{equation*}
  \frac{d}{dx} f(x) = x + 2 x \ln\left(x\right)
\end{equation*}
\begin{equation*}
  f'(1) = 1
\end{equation*}
A second derivative:

\begin{equation*}
  \frac{d^2}{dx^2} \sin(3 x) = -9 \sin\left(3 x\right)
\end{equation*}
A partial derivative:

\begin{equation*}
  \frac{\partial}{\partial y} (x^2 y^3) = 3 x^{2} y^{2}
\end{equation*}
\end{document}
